\documentclass{article}
\usepackage[utf8]{inputenc}
\usepackage{amsmath,amssymb}
\usepackage[most]{tcolorbox}
\usepackage{geometry}
\geometry{margin=2.5cm}

\newcommand{\step}[1]{\par\noindent\hangindent=3.5em\hangafter=1%
  \makebox[3.5em][l]{\textbf{#1.}}}
\newcommand{\steptag}[1]{\hfill{\small\textsf{[#1]}}}

\newtcolorbox{proofbox}[1]{
  colback=gray!5, colframe=gray!60,
  fonttitle=\bfseries, title={#1},
  breakable, enhanced,
  left=4pt, right=4pt, top=4pt, bottom=4pt,
  before skip=10pt, after skip=10pt
}

\begin{document}

\begin{proofbox}{For unital positive Phi with ||Phi\textasciicircum{}2-Phi|| <= eta <= eta\_...}
\step{1} For unital positive Phi with ||Phi\textasciicircum{}2-Phi|| <= eta <= eta\_0, P=theta(2Phi-1), A=Im P, if Phi admits a faithful state omega with ||omega∘Phi-omega|| <= eta and density rho\_omega >= lambda 1 (lambda>0), then for all a,b in A the ambient product hole satisfies ||a∘b-P(a∘b)|| = ||h\_\{a,b\}|| <= C (eta/lambda) ||a|| ||b||, with universal dimension-free C. \steptag{validated} \\[6pt]
\step{1.1} SPECTRAL SPLIT + q\textasciicircum{}- BOUND. q\_a=P(a\textasciicircum{}2)-a\textasciicircum{}2 is self-adjoint in Ker P (GT-hyp). By functional calculus (node 1.1.1) split q\_a=q\textasciicircum{}+ - q\textasciicircum{}- with q\textasciicircum{}+,q\textasciicircum{}- >= 0, q\textasciicircum{}+ o q\textasciicircum{}- = 0, and ||q\_a||=max(||q\textasciicircum{}+||,||q\textasciicircum{}-||). By GT-squarehole (instantiate r:=a): q\_a >= -C eta ||a||\textasciicircum{}2 1. Hence 0 <= q\textasciicircum{}- <= C eta ||a||\textasciicircum{}2 1 (GT-orderunit-def: -c1<=q\_a means q\textasciicircum{}-<=c1), so ||q\textasciicircum{}-|| <= C eta ||a||\textasciicircum{}2 (GT-sup-states monotonicity, 0<=q\textasciicircum{}-<=C eta||a||\textasciicircum{}2 1). \steptag{validated} \\[6pt]
\step{1.1.1} FUNCTIONAL CALCULUS SPLIT. V=B(H)\_sa is a unital JB algebra. For self-adjoint q=q\_a, by GT-jb-spectral C(q) is isometrically Jordan-isomorphic (phi) to C(Sp q), phi(q)=iota. Set q\textasciicircum{}+ = phi\textasciicircum{}\{-1\}(t\textasciicircum{}+), q\textasciicircum{}- = phi\textasciicircum{}\{-1\}(t\textasciicircum{}-) where t\textasciicircum{}+=max(t,0), t\textasciicircum{}-=max(-t,0) on Sp q. Then q\textasciicircum{}+,q\textasciicircum{}- in C(q), q\textasciicircum{}+,q\textasciicircum{}- >= 0 (GT-jb-positive-spectrum: pointwise nonneg => positive), q\textasciicircum{}+ - q\textasciicircum{}- = phi\textasciicircum{}\{-1\}(t)=q, and q\textasciicircum{}+ o q\textasciicircum{}- = phi\textasciicircum{}\{-1\}(t\textasciicircum{}+ t\textasciicircum{}-)=0 (t\textasciicircum{}+ t\textasciicircum{}- =0 pointwise). Sp q = Sp(q\textasciicircum{}+) cup (-Sp(q\textasciicircum{}-)) (disjoint supports), so by GT-norm-spectral ||q||=sup|Sp q|=max(sup Sp q\textasciicircum{}+, sup Sp q\textasciicircum{}-)=max(||q\textasciicircum{}+||,||q\textasciicircum{}-||) (q\textasciicircum{}+,q\textasciicircum{}->=0 so their norms equal their max eigenvalues). \steptag{validated} \\[4pt]
\step{1.2} EXPECTATION BOUND |omega(q\_a)| <= C eta ||a||\textasciicircum{}2. omega(q\_a)=omega(P(a\textasciicircum{}2))-omega(a\textasciicircum{}2). Write omega(P(a\textasciicircum{}2))=omega(Phi(a\textasciicircum{}2)) + omega((P-Phi)(a\textasciicircum{}2)). By GT-Pprops delta=||P-Phi||<=C eta; omega is a state so |omega(y)|<=||y|| (GT-state-bounded); ||a\textasciicircum{}2||=||a||\textasciicircum{}2 (GT-bh-cstar). So |omega((P-Phi)(a\textasciicircum{}2))| <= ||(P-Phi)(a\textasciicircum{}2)|| <= delta||a\textasciicircum{}2|| <= C eta||a||\textasciicircum{}2. Also omega(Phi(a\textasciicircum{}2))=(omega o Phi)(a\textasciicircum{}2)=omega(a\textasciicircum{}2)+((omega o Phi)-omega)(a\textasciicircum{}2), and |((omega o Phi)-omega)(a\textasciicircum{}2)| <= ||omega o Phi - omega||*||a\textasciicircum{}2|| <= eta||a||\textasciicircum{}2 (GT-hyp). Hence omega(P(a\textasciicircum{}2))=omega(a\textasciicircum{}2)+r with |r|<=C eta||a||\textasciicircum{}2, so |omega(q\_a)|=|r|<=C eta||a||\textasciicircum{}2. \steptag{validated} \\[4pt]
\step{1.3} omega(q\textasciicircum{}+) <= C eta ||a||\textasciicircum{}2. Since q\_a=q\textasciicircum{}+-q\textasciicircum{}- (1.1), omega(q\textasciicircum{}+)=omega(q\_a)+omega(q\textasciicircum{}-). |omega(q\_a)|<=C eta||a||\textasciicircum{}2 (1.2). q\textasciicircum{}->=0 and omega a state, so 0<=omega(q\textasciicircum{}-)<=||q\textasciicircum{}-|| (GT-state-bounded / GT-sup-states) <= C eta||a||\textasciicircum{}2 (1.1). Hence omega(q\textasciicircum{}+)<=C eta||a||\textasciicircum{}2. \steptag{validated} \\[4pt]
\step{1.4} FAITHFULNESS UPGRADE ||q\textasciicircum{}+|| <= C (eta/lambda) ||a||\textasciicircum{}2. For any x>=0: omega(x)=Tr(rho\_omega x) (GT-state-density). By node 1.4.1, Tr(rho\_omega x) >= lambda Tr(x); by node 1.4.2, Tr(x) >= ||x||. So omega(x) >= lambda||x|| for x>=0. Apply to x=q\textasciicircum{}+ >= 0 (the positive part from node 1.1, a transitive dependency of this node via 1.3 whose declared deps are [1.1,1.2]; 1.1 is thus a recorded validated ancestor of 1.4): lambda||q\textasciicircum{}+|| <= omega(q\textasciicircum{}+) <= C eta||a||\textasciicircum{}2 (node 1.3), so ||q\textasciicircum{}+|| <= C (eta/lambda)||a||\textasciicircum{}2. \steptag{validated} \\[6pt]
\step{1.4.1} Tr(rho\_omega x) >= lambda Tr(x) for x>=0. By GT-hyp rho\_omega>=lambda1, so rho\_omega-lambda1>=0 (GT-orderunit-def). By GT-trace-psd (self-duality of the positive cone of L(H): a in J\textasciicircum{}+ iff tr(ab)>=0 for all b in J\textasciicircum{}+), with a:=rho\_omega-lambda1>=0 and b:=x>=0, Tr((rho\_omega-lambda1)x)>=0, i.e. Tr(rho\_omega x) >= lambda Tr(x). \steptag{validated} \\[4pt]
\step{1.4.2} Tr(x) >= ||x|| for x>=0. By GT-eigen-nonneg, x>=0 means x self-adjoint with all eigenvalues mu\_i>=0. By GT-norm-spectral ||x||=sup|Sp x|=max\_i mu\_i (mu\_i>=0). Tr(x)=sum\_i mu\_i >= max\_i mu\_i = ||x|| (all summands nonneg). \steptag{validated} \\[4pt]
\step{1.5} DIAGONAL ||q\_a|| <= C (eta/lambda)||a||\textasciicircum{}2. ||q\_a||=max(||q\textasciicircum{}+||,||q\textasciicircum{}-||) (1.1). ||q\textasciicircum{}+||<=C(eta/lambda)||a||\textasciicircum{}2 (1.4). ||q\textasciicircum{}-||<=C eta||a||\textasciicircum{}2 (1.1) <= C(eta/lambda)||a||\textasciicircum{}2 since lambda<=1 (GT-hyp). Hence ||q\_a||<=C(eta/lambda)||a||\textasciicircum{}2. \steptag{validated} \\[4pt]
\step{1.6} POLARISATION => the contract. h\_\{a,b\}=a∘b-P(a∘b)=-q\_\{a,b\} (GT-hyp/def-square-hole) is symmetric bilinear in (a,b), with diagonal h\_\{a,a\}=-q\_a, so ||h\_\{a,a\}||=||q\_a||<=C(eta/lambda)||a||\textasciicircum{}2 (1.5). For t>0, h\_\{a,b\}=(1/(4t))(h\_\{a+tb,a+tb\}-h\_\{a-tb,a-tb\}). Hence ||h\_\{a,b\}|| <= (1/(4t))(||q\_\{a+tb\}||+||q\_\{a-tb\}||) <= (C eta/(4 t lambda))(||a+tb||\textasciicircum{}2+||a-tb||\textasciicircum{}2) <= (C eta/(2 t lambda))(||a||+t||b||)\textasciicircum{}2 (triangle inequality on the norm). For a,b!=0 choose t=||a||/||b||: (||a||+t||b||)\textasciicircum{}2=4||a||\textasciicircum{}2 and 1/(2t)=||b||/(2||a||), giving ||h\_\{a,b\}|| <= C(eta/lambda)||a|| ||b|| (the zero cases are trivial). This is the contract. \steptag{validated} \\[4pt]
\end{proofbox}

\end{document}
